\subsection{标量环的更换}

设 A、B 为两个环，且 \(\rho\) 为从环B到环A的一个同态。对于每个A-模E，外合成律 \((\beta, x) \mapsto \rho(\beta)x\) （连同加法）定义了一个B-模结构，称为与 \(\rho\) 及E上的A-模结构相关联；这个B-模记为 \(\rho_{*}(\mathbf{E})\) 或 \(E_{[B]}\) （若无混淆，甚至简记为E）。特别地，若B是A的子环且 \(\rho: B \to A\) 是典范单射， \(E_{[B]}\) 称为由将标量环A限制到B而得到的B-模；出于语言上的滥用，当同态 \(\rho\) 是任意的时也使用这一表述。

若 \(\mathbf{F}\) 是A-模 \(\mathbf{E},\rho_{*}(\mathbf{F})\) 是 \(\rho_{*}(\mathbf{E})\) 的一个子模，且 \(\rho_{*}(\mathbf{E} / \mathbf{F})\) 等于 \(\rho_{*}(\mathbf{E}) / \rho_{*}(\mathbf{F})\) .

设E、F为两个A-模；每个A-线性映射 \(u: E \rightarrow F\) 也是B-线性映射 \(E_{[B]} \rightarrow F_{[B]}\) ，记为 \(\rho_{*}(u)\) ；换言之，存在Z-模的一个典范单射

\[
\operatorname{Hom} _ {\mathbf {A}} (\mathbf {E}, \mathbf {F}) \rightarrow \operatorname{Hom} _ {\mathbf {B}} (\mathbf {E} _ {[ \mathbf {B} ]}, \mathbf {F} _ {[ \mathbf {B} ]}).\tag{47}
\]

此映射未必是双射；换言之，B-线性映射 \(E_{[B]} \rightarrow F_{[B]}\) 未必是A-线性的。例如， \(E_{[B]}\) 的一个子B-模未必是E的子A-模：若A是域且B是A的子域，则 \(B_{s}\) 是向量B-空间 \((A_{s})_{[B]}\) 的一个向量子空间，但若 \(B \neq A\) ，它未必是向量子A-空间.

显然，对于每一个A-模族 \((\mathbf{E}_t)_{t\in I}\) ，B-模 \(\rho_{*}\left(\prod_{i\in I}\mathbf{E}_{i}\right)\left(\text{resp. } \rho_{*}\left(\bigoplus_{i\in I}\mathbf{E}_{i}\right)\right)\) 等于 \(\prod_{i\in I}\rho_{*}(\mathbf{E}_{i})\) （相应地， \(\bigoplus_{i\in I}\rho_{*}(\mathbf{E}_{i})\) ).

\(\rho_{*}(\mathbf{E})\) 的每个生成系都是 E 的一个生成系，但反之不一定成立。

命题 24. 设 A, B 是两个环，且 \(\rho: \mathbf{B} \to \mathbf{A}\) 一个环同态。

\begin{enumerate}
\def\labelenumi{(\roman{enumi})}
\item
  若 \(\rho\) 是满射，则典范映射（47）是双射。对于每个A-模E， \(\rho_{*}(\mathbf{E})\) 的每个子B-模都是 E 的一个子A-模；E 的每个生成系都是 \(\rho_{*}(\mathbf{E})\) 的生成系.
\item
  若 \(\rho\) 是单射的，则A-模 \(\mathbf{E}\) 中的每个自由族都是B-模 \(\rho_{*}(\mathbf{E})\) 中的自由族.
\end{enumerate}

该命题由定义直接得出。

请注意，即使 \(\rho\) 是单射的， \(\rho_{*}(\mathbf{E})\) 中的自由族在 \(\mathbf{E}\) 中也未必是自由的.

*例如，1 和 \(\sqrt{2}\) 在把 \(\mathbf{R}\) 视为向量 \(\mathbf{R}\) -空间时并不构成自由系，尽管在把 \(\mathbf{R}\) 视为向量 \(\mathbf{Q}\) -空间时它们构成一个自由系（参见注记1）。*

命题25. 设A、B为两个环， \(\rho: \mathbf{B} \to \mathbf{A}\) 一个环同态，且E是一个A-模。设 \((\alpha_{\lambda})_{\lambda \in L}\) 是A视为左B-模时的一个生成系（相应地，元素的自由族、基）。设 \((a_{\mu})_{\mu \in M}\) 是A-模E的一个生成系（相应地，自由元素族、基）。则 \((\alpha_{\lambda}a_{\mu})_{(\lambda,\mu) \in L \times M}\) 是一个生成系（相应地，（当 \(\rho\) 是单射时）是自由元素族、基），都是就B-模 \(\rho_{*}(\mathbf{E})\) 而言.

若 \(x = \sum_{\mu \in \mathbf{M}} \gamma_{\mu} a_{\mu}\) ，其中 \(\gamma_{\mu} \in \mathbf{A}\) ，而 \((\alpha_{\lambda})\) 是 \(\mathbf{A}\) 的一个生成系，则可写 \(\gamma_{\mu} = \sum_{\lambda \in L} \rho(\beta_{\lambda \mu}) \alpha_{\lambda}\) ，其中 \(\beta_{\lambda \mu} \in \mathbf{B}\) ，对所有 \(\mu \in \mathbf{M}\) 成立，由此 \(x = \sum_{\lambda, \mu} \rho(\beta_{\lambda \mu}) \alpha_{\lambda} a_{\mu}\) 。另一方面，若 \((\alpha_{\lambda})\) 与 \((a_{\mu})\) 是自由族，则关系式 \(\sum_{\lambda, \mu} \rho(\beta_{\lambda \mu}) \alpha_{\lambda} a_{\mu} = 0\) （其中 \(\beta_{\lambda \mu} \in \mathbf{B}\) ）可写成 \(\sum_{\mu \in M} \left( \sum_{\lambda \in L} \rho(\beta_{\lambda \mu}) \alpha_{\lambda} \right) a_{\mu} = 0\) ；因此它蕴含 \(\sum_{\lambda \in L} \rho(\beta_{\lambda \mu}) \alpha_{\lambda} = 0\) 对所有 \(\mu \in M\) 成立，从而 \(\beta_{\lambda \mu} = 0\) 对所有 \(\lambda, \mu\) 成立，只要 \(\rho\) 是单射。

推论. 若 A 是有限生成的左 B-模，且 E 是有限生成的左 A-模，则 \(\rho_{*}(\mathbf{E})\) 是有限生成的左 A-模。

设 C 为第三个环， \(\rho':\mathbf{C}\to\mathbf{B}\) 一个环同态，且 \(\rho''=\rho\circ\rho'\) 为复合同态。由定义立即得出，对每个A-模E有 \(\rho_{*}^{\prime \prime}(\mathbf{E}) = \rho_{*}^{\prime}(\rho_{*}(\mathbf{E}))\) 。特别地，若 \(\rho\) 是B到A上的一个同构，则 \(\mathbf{E} = \rho_{*}^{\prime}(\rho_{*}(\mathbf{E}))\) ，其中 \(\rho^{\prime}\) 表示 \(\rho\) 的逆同构.

注：(1) 设K是一个域，A是K的一个子环，且具有如下性质：对K的元素的每个有限族 \((\xi_{i})_{1\leqslant i\leqslant n}\) ，存在一个非零的 \(\gamma\in A\) ，使得 \(\gamma\xi_{i}\in A\) 对于 \(1\leqslant i\leqslant n\) 成立（当A是交换的且K是A的分式域时，这一假设总是满足的）。设E是K上的向量空间， \(E_{[A]}\) 是由将标量环限制到A而得到的A-模。那么，若族 \((x_{\lambda})_{\lambda\in L}\) 在 \(E_{[A]}\) 中是自由的，则它在E中也是自由的。注意可局限于 \(L=\{1,n\}\) 的情形；若存在关系 \(\sum_{i=1}^{n} \xi_i x_i = 0\) ，其中 \(\xi_i \in \mathbf{K}\) ， \(\xi_i\) 并非全为零，则对所有 \(\beta \in \mathbf{A}, \sum_{i=1}^{n} (\beta \xi_i) x_i = 0\) 。根据假设，可以假定 \(\beta \neq 0\) 在 \(\mathbf{A}\) 中，使得 \(\beta \xi_{i} = \alpha_{i}\) 对所有 \(i\) 都属于A；但关系 \(\sum_{i=1}^{n} \alpha_{i} x_{i} = 0\) 与假设相矛盾，因为 \(\alpha_{i}\) 并非全为零。

\begin{enumerate}
\def\labelenumi{(\arabic{enumi})}
\setcounter{enumi}{1}
\tightlist
\item
  如果环同态 \(\rho: \mathbf{B} \to \mathbf{A}\) 是满射且 \(\mathfrak{b}\) 是其核（因此 \(\mathbf{A}\) 被典范地等同于 \(\mathbf{B} / \mathfrak{b}\) ），则对于每个 A-模 \(\mathbf{E}\) , \(\mathfrak{b}\) 包含于 \(\rho_*(\mathbf{E})\) 的零化子中，且 \(\mathbf{E}\) 是由 \(\rho_*(\mathbf{E})\) 按第12号中定义的过程得到的A-模。
\end{enumerate}

设 A、B 为两个环，且 \(\rho: \mathbf{B} \to \mathbf{A}\) 一个同态。设E是一个A-模，F是一个B-模；B-线性映射 \(u: \mathbf{F} \to \rho_{*}(\mathbf{E})\) （若无混淆，也称为从F到E的B-线性映射）也称为（相对于 \(\rho\) ）把B-模F映入A-模E的半线性映射；也称有序对 \((u, \rho)\) 为F到E的一个双态射；因此这意味着，对于 \(x \in \mathbf{F}\) , \(y \in \mathbf{F}\) 和 \(\beta \in \mathbf{B}\) ,

\[
\left\{ \begin{array}{l} u (x + y) = u (x) + u (y) \\ u (\beta x) = \rho (\beta) u (x). \end{array} \right.\tag{48}
\]

集合 \(\operatorname{Hom}_{\mathbf{B}}(\mathbf{F},\rho_{*}(\mathbf{E}))\) 由把 \(\mathbf{F}\) 映入 \(\mathbf{E}\) 的B-线性映射组成，若不会引起混淆，也写作 \(\operatorname{Hom}_{\mathbf{B}}(\mathbf{F},\mathbf{E})\) 。

当 \(\rho\) 是B到A上的同构时，关系 \(u(\beta x) = \rho(\beta) u(x)\) （对所有 \(\beta \in B\) ）也可写成 \(u(\rho'(\alpha)x) = \alpha x\) （对所有 \(\alpha \in A\) ），其中 \(\rho'\) 表示 \(\rho\) 的逆同构；说 \(u\) 相对于 \(\rho\) 是半线性的，于是等价于说 \(u\) 是 \(\rho_*(F)\) 到 E 的一个A-线性映射。

例. 已经看到（第1号），位似 \(h_{\alpha} \colon x \mapsto \alpha x\) 在A-模 \(\mathbf{E}\) 上不一定是线性映射。但如果 \(\alpha\) 是可逆的，则 \(h_{\alpha}\) 是相对于内自同构 \(\xi \mapsto \alpha \xi \alpha^{-1}\) （ \(\mathbf{A}\) 的内自同构）的半线性映射（而且还是双射的），因为 \(\alpha(\lambda x) = (\alpha \lambda \alpha^{-1})(\alpha x)\) .

设 C 为第三个环， \(\rho':\mathrm{C}\to\mathrm{B}\) 一个同态，且 G 是一个 C-模。如果 \(v:\mathrm{G}\to\mathrm{F}\) 是相对于 \(\rho'\) 的半线性映射，则复合 \(w=u\circ v\) 是 G 到 E 的相对于同态 \(\rho''=\rho\circ\rho'\) 的半线性映射。若 \(\rho\) 是一个同构且 \(u:\mathrm{F}\to\mathrm{E}\) 是相对于 \(\rho\) 的双射半线性映射，则其逆映射 \(u':\mathrm{E}\to\mathrm{F}\) 是相对于逆同构 \(\rho':\mathrm{A}\to\mathrm{B}\) （ \(\rho\) 的逆同构）的半线性映射.

由此可见，对于通过在集合的有序对 (A, E) 上给定 A 上的环结构和 E 上的左 A-模结构所定义的结构种，其双态射 \((u, \phi)\) 可取作态射（集合论，IV，§2，第1号）；在以下内容中，我们始终假定已作出这一态射选择。

注记（3）. 设 \(A_{1}\) , \(A_{2}\) 为两个环， \(A = A_{1} \times A_{2}\) 为它们的积，并令 \(e_{1} = (1, 0)\) , \(e_{2} = (0, 1)\) 在 A 中，使得 \(A_{1}\) 与 \(A_{2}\) 分别典范地等同于双边理想 \(Ae_{1}\) 与 \(Ae_{2}\) 。对于每个 A-模 E， \(e_{1}E\) 和 \(e_{2}E\) 是E的子A-模 \(E_{1}\) , \(E_{2}\) ，分别被 \(e_{2}\) 和 \(e_{1}\) 零化，因此，在把 \(A/Ae_{2}\) 与 \(A_{1}\) 、 \(A/Ae_{1}\) 与 \(A_{2}\) 典范地等同之后， \(E_{1}\) （相应地 \(E_{2}\) ）被赋予一个 \(A_{1}\) -模（相应地， \(A_{2}\) -模）结构。此外，E 是 \(E_{1}\) 和 \(E_{2}\) 的直和，因为每一个 \(x \in E\) 都可以写成 \(x = e_{1}x + e_{2}x\) ，以及关系 \(e_{1}x = e_{2}y\) 蕴含 \(e_{1}x = e_{1}^{2}x = e_{1}e_{2}y = 0\) . 反之，对于由一个 \(A_{1}\) -模 \(F_{1}\) 和一个 \(A_{2}\) -模 \(F_{2}\) 组成的每个有序对，设 \(E_{1}\) 为A-模 \((p_{1}) * (F_{1})\) , \(E_{2}\) 为A-模 \((p_{2}) * (F_{2})\) ，其中 \(p_{1}\) 和 \(p_{2}\) 是A到 \(A_{1}\) 和 \(A_{2}\) 上的投影；则在A-模 \(E = E_{1} \oplus E_{2}\) 中， \(E_{1} = e_{1}E\) , \(E_{2} = e_{2}E\) 。因此，对A-模的研究归结为对 \(A_{1}\) -模与 \(A_{2}\) -模的研究。特别地，E 的每个子模 M 都具有形式 \(M_{1} \oplus M_{2}\) ，其中 \(M_{1} = e_{1}M\) 和 \(M_{2} = e_{2}M\) .

